\section{Moral outrage: this should not happen in America}
\label{sec:outrage}

Many health laws gain momentum after a perceived injustice: an insurance denial,
an unaffordable treatment, a discriminatory algorithm, or plainly unequal access
to care. The framing becomes a sentence a member can repeat on the floor. This is
unfair, and it should not happen here. Outrage is effective precisely because it
draws both public attention and press coverage to a specific, fixable wrong, which
is why coalition building around a documented injustice is such a powerful
legislative move.

The injustice here is documented, not asserted. A widely used clinical risk
algorithm was shown to systematically underestimate the needs of Black patients,
because it used past spending as a proxy for illness, and spending itself
reflected unequal access \cite{obermeyer2019dissecting}. An unverified algorithm
can encode a disparity and then scale it across every patient it touches. The
author's comparison of proposed bills places patient priority, rather than
institutional convenience, at the center of the remedy
\cite{Kawchak2026PatientPriorityProposedU}. \autoref{fig:outrage} shows how a
documented gap becomes an audited safeguard.

\begin{narrativefig}
\node[nfrisk] (a) at (0,0) {Unequal access};
\node[nfrisk] (b) at (3.5,0) {Bias encoded};
\node[nfdec] (c) at (7.0,0) {Bias surveillance};
\node[nfact] (d) at (10.6,0) {Measured and reported};
\node[nfhope] (e) at (10.6,-2.1) {Equal protection};
\draw[nfedge] (a)--(b); \draw[nfedge] (b)--(c); \draw[nfedge] (c)--(d);
\draw[nfedge] (d)--(e);
\end{narrativefig}
\vspace{-0.15cm}
\figcaption{Figure 5. From a documented disparity to an audited safeguard
(Mermaid 15 as colored TikZ). Surveillance turns an unfair gap into a measured,
correctable number.}

\autoref{tab:outrage} contrasts the unverified status quo with the remedy the bill
provides, so the outrage points somewhere constructive.

\begin{table}[H]
\footnotesize
\begin{tabularx}{\textwidth}{@{}L{4.2cm} Y Y@{}}
\toprule
\textbf{Concern} & \textbf{Unverified status quo} & \textbf{Remedy in H. R. 9510}\\
\midrule
Algorithmic bias & Hidden, scaled silently & Subgroup surveillance, reported\\
Unequal access & Tolerated as background & Access recorded and audited\\
Accountability & Diffuse, hard to trace & Named record, public report\\
Correction & Slow or never & Triggered by a measured gap\\
\bottomrule
\end{tabularx}
\caption{Table 3. Turning a documented injustice into a statutory check.}
\label{tab:outrage}
\end{table}

The remedy is not outrage for its own sake. It is a precise statutory check that
converts a documented disparity into a measured, reported, and correctable number.
A committee that has seen the data can ask the right question. If we can measure
the unfairness, why would we choose not to?
